\documentclass[11pt]{article}

\usepackage[margin=1in]{geometry}
\usepackage{amsmath,amssymb,amsthm,mathtools}
\usepackage[T1]{fontenc}
\usepackage[utf8]{inputenc}
\usepackage{lmodern}
\usepackage{hyperref}
\usepackage{enumitem}

\newtheorem{theorem}{Theorem}
\newtheorem{lemma}{Lemma}
\newtheorem{definition}{Definition}
\newtheorem{remark}{Remark}
\newtheorem{proposition}{Proposition}

\title{Deterministic Closure Obstruction via CFI-Type Local--Global Decoupling}
\author{Inacio F. Vasquez}
\date{\today}

\begin{document}
\maketitle

\begin{abstract}
We record the deterministic-closure obstruction in a corrected formal form. For every fixed locality radius \(R\) and variable bound \(k\), we construct an explicit infinite family of bounded-degree graph pairs \((G_0^{(n)},G_1^{(n)})\) with identical rooted \(R\)-ball multisets and hence identical deterministic \(\mathrm{FO}^k\)-local transcript data, but with distinct global parity invariant. The witness family is given by CFI-type parity twists over high-girth bounded-degree base graphs. This yields a local--global decoupling theorem: no deterministic \(\mathrm{FO}^k\)-local refinement rule can recover the global invariant from the local transcript. Therefore deterministic closure fails in the refinement-only regime.
\end{abstract}

\section{Setup}

Fix integers \(R\ge 1\) and \(k\ge 1\). Let \(\mathcal G\) denote the class of finite graphs.

\begin{definition}[Rooted \(R\)-ball multiset]
For a graph \(G\) and \(v\in V(G)\), let \(B_R^G(v)\) denote the rooted radius-\(R\) neighborhood of \(v\). Define
\[
T_R(G):=\big\{\,[B_R^G(v)]_{\cong}\ :\ v\in V(G)\,\big\},
\]
the multiset of rooted isomorphism classes of radius-\(R\) balls.
\end{definition}

\begin{definition}[Deterministic \(\mathrm{FO}^k\)-local refinement]
A deterministic \(\mathrm{FO}^k\)-local refinement rule at locality radius \(R\) is any map
\[
\mathcal A:\mathcal G\to \mathcal O
\]
whose output depends only on the rooted radius-\(R\) local transcript. Equivalently, there exists a function
\[
f_R
\]
such that
\[
\mathcal A(G)=f_R(T_R(G))
\]
for all \(G\in\mathcal G\).
\end{definition}

This is the only dependence property needed for the obstruction argument.

\section{Base graphs and CFI parity gadgets}

Let \(B\) be a connected \(d\)-regular graph with \(d\ge 3\) and girth
\[
\mathrm{girth}(B)>2R.
\]
We will later take an infinite family \(B_n\) with uniformly bounded degree and girth tending to infinity.

For each vertex \(v\in V(B)\), let \(E(v)\) denote the set of incident edges. For \(i\in\{0,1\}\), define the vertex set
\[
V(G_i):=\{(v,\alpha)\ :\ v\in V(B),\ \alpha:E(v)\to \mathbb F_2,\ \sum_{e\in E(v)}\alpha(e)=i\}.
\]
Thus each base vertex \(v\) is replaced by the parity class of all local edge-labelings of parity \(i\).

For each base edge \(e=uv\in E(B)\), add an edge in \(G_i\) between \((u,\beta)\) and \((v,\alpha)\) whenever
\[
\beta(e)=\alpha(e).
\]
This is the standard parity-consistent CFI gluing rule.

\begin{remark}
The graph \(G_0\) is the even-parity lift and \(G_1\) is the odd-parity lift. Both have bounded degree depending only on \(d\), and
\[
|V(G_i)|=|V(B)|\,2^{d-1}.
\]
\end{remark}

\section{Local coincidence}

\begin{lemma}[Radius-\(R\) rooted balls coincide]
If \(\mathrm{girth}(B)>2R\), then
\[
T_R(G_0)=T_R(G_1).
\]
\end{lemma}

\begin{proof}
Fix \(v\in V(B)\) and vertices \((v,\alpha_0)\in V(G_0)\), \((v,\alpha_1)\in V(G_1)\). Because \(\mathrm{girth}(B)>2R\), the radius-\(R\) neighborhood of \(v\) in \(B\) is a tree. Hence every lift of this neighborhood to either \(G_0\) or \(G_1\) is determined recursively from the root choice, with no cycle constraints encountered inside depth \(R\).

At each step away from the root, a child choice is determined by one shared edge-bit and free choices on the remaining \(d-1\) incident edges subject to one parity equation. Therefore the rooted branching pattern is identical in \(G_0\) and \(G_1\). Since there are no cycles within radius \(R\), no contradiction can arise from transporting the local labeling data.

Thus
\[
B_R^{G_0}(v,\alpha_0)\cong B_R^{G_1}(v,\alpha_1)
\]
as rooted graphs. Since this holds at every root, the multisets of rooted \(R\)-ball isomorphism types coincide:
\[
T_R(G_0)=T_R(G_1).
\]
\end{proof}

\section{Global invariant}

Fix a spanning tree \(T\subseteq B\). For any realization \(G\) of the parity construction above, define the twist bit on a non-tree edge \(e=uv\in E(B)\setminus E(T)\) by transporting the local labels from a chosen root along the unique tree paths in \(T\) to \(u\) and \(v\), then comparing the edge-bit assigned to \(e\) at the two endpoints. This yields a bit
\[
\tau_G(e)\in \mathbb F_2.
\]
Define the parity invariant
\[
I(G):=\sum_{e\in E(B)\setminus E(T)} \tau_G(e)\pmod 2.
\]

\begin{lemma}[Well-definedness]
The invariant \(I(G)\) is independent of the root choice and of the transported representatives.
\end{lemma}

\begin{proof}
Changing the root or transported representatives modifies the local edge-bit system by a coboundary on the base graph \(B\). The induced change on the vector \((\tau_G(e))_{e\notin T}\) is addition of a cut-space element. Summing over the fundamental cycle coordinates modulo \(2\) is invariant under such coboundary changes. Hence \(I(G)\) is well-defined.
\end{proof}

\begin{lemma}[Parity separation]
For the two parity lifts \(G_0,G_1\),
\[
I(G_0)\neq I(G_1).
\]
\end{lemma}

\begin{proof}
By construction, the even lift \(G_0\) has trivial total parity class, while the odd lift \(G_1\) differs from it by the unique nontrivial global parity twist. In the spanning-tree coordinate system, this changes the total parity on the fundamental cycle coordinates by \(1\in\mathbb F_2\). Therefore
\[
I(G_0)=0,\qquad I(G_1)=1.
\]
Hence \(I(G_0)\neq I(G_1)\).
\end{proof}

\section{\(\mathrm{FO}^k\)-local indistinguishability}

\begin{lemma}[Local logic cannot separate the pair]
For fixed \(R,k\), if \(\mathrm{girth}(B)>2R\), then
\[
G_0\equiv_{\mathrm{FO}^k_R} G_1.
\]
\end{lemma}

\begin{proof}
A formula in \(\mathrm{FO}^k\) of locality radius \(R\) only inspects rooted radius-\(R\) neighborhoods and the incidence relations among finitely many pebbled points inside such neighborhoods. Since all rooted radius-\(R\) neighborhoods agree between \(G_0\) and \(G_1\), the duplicator can respond in the \(k\)-pebble \(R\)-round Ehrenfeucht--Fra\"\i ss\'e game by matching the corresponding rooted tree pattern at each move.

Because the base graph has girth \(>2R\), these neighborhoods are tree-like and contain no conflicting cycle witness visible within radius \(R\). Thus the duplicator preserves partial isomorphism for all \(R\) rounds. Therefore
\[
G_0\equiv_{\mathrm{FO}^k_R} G_1.
\]
\end{proof}

\section{Deterministic closure obstruction}

\begin{theorem}[Deterministic closure fails]
For every fixed \(R,k\), there exists an explicit infinite family of bounded-degree graph pairs
\[
\{(G_0^{(n)},G_1^{(n)})\}_{n\ge 1}
\]
such that
\[
T_R(G_0^{(n)})=T_R(G_1^{(n)}),
\]
\[
G_0^{(n)}\equiv_{\mathrm{FO}^k_R} G_1^{(n)},
\]
\[
I(G_0^{(n)})\neq I(G_1^{(n)}).
\]
Consequently, no deterministic \(\mathrm{FO}^k\)-local refinement rule can recover the invariant \(I\) from the local transcript.
\end{theorem}

\begin{proof}
Choose any explicit infinite family \((B_n)\) of bounded-degree graphs with girth tending to infinity, for example any explicit \(d\)-regular high-girth family. For \(n\) large enough that
\[
\mathrm{girth}(B_n)>2R,
\]
construct the corresponding parity lifts \(G_0^{(n)},G_1^{(n)}\).

By the previous lemmas,
\[
T_R(G_0^{(n)})=T_R(G_1^{(n)}),
\qquad
G_0^{(n)}\equiv_{\mathrm{FO}^k_R} G_1^{(n)},
\qquad
I(G_0^{(n)})\neq I(G_1^{(n)}).
\]

Now let \(\mathcal A\) be any deterministic \(\mathrm{FO}^k\)-local refinement rule. By definition there exists a function \(f_R\) such that
\[
\mathcal A(G)=f_R(T_R(G)).
\]
Hence
\[
\mathcal A(G_0^{(n)})=f_R(T_R(G_0^{(n)}))=f_R(T_R(G_1^{(n)}))=\mathcal A(G_1^{(n)}).
\]
Therefore \(\mathcal A\) cannot distinguish the pair, while the global invariant does distinguish the pair:
\[
I(G_0^{(n)})\neq I(G_1^{(n)}).
\]
So \(I\) is not determined by the deterministic local transcript. Deterministic closure fails.
\end{proof}

\section{Closed-form consequence}

\begin{corollary}[Local--global decoupling]
For every fixed \(R\),
\[
I(G)\notin \sigma(T_R(G))
\]
on the witness family. Equivalently,
\[
\not\exists f:\operatorname{Im}(T_R)\to \mathbb F_2
\quad\text{such that}\quad
f(T_R(G))=I(G)
\]
for all graphs \(G\) in the witness family.
\end{corollary}

\begin{proof}
If such an \(f\) existed, then identical \(T_R\)-data would force identical \(I\)-values. This contradicts
\[
T_R(G_0^{(n)})=T_R(G_1^{(n)})
\quad\text{and}\quad
I(G_0^{(n)})\neq I(G_1^{(n)}).
\]
\end{proof}

\section{Status statement}

The deterministic closure claim in the local refinement regime is therefore false in full generality. The obstruction is explicit, bounded-degree, infinite, and CFI-realized.

\[
\boxed{\text{Deterministic closure fails for local refinement systems.}}
\]

\end{document}
